\documentclass[conference]{IEEEtran}
\IEEEoverridecommandlockouts

% Configuración para español
\usepackage[spanish,es-nodecimaldot,es-noshorthands]{babel}
\usepackage[utf8]{inputenc}
\usepackage[T1]{fontenc}

% Paquetes estándar IEEE
\usepackage{cite}
\usepackage{amsmath,amssymb,amsfonts}
\usepackage{algorithmic}
\usepackage{graphicx}
\usepackage{textcomp}
\usepackage{xcolor}
\usepackage{float}
\usepackage{listings} % Para código

% ---------- Configuración de listados de código ----------
\definecolor{codegray}{rgb}{0.45,0.45,0.45}
\definecolor{codeblue}{rgb}{0.10,0.10,0.60}
\lstset{
	basicstyle=\ttfamily\scriptsize,
	keywordstyle=\color{codeblue}\bfseries,
	commentstyle=\color{codegray}\itshape,
	stringstyle=\color{codegray},
	numbers=left,
	numberstyle=\tiny\color{codegray},
	numbersep=5pt,
	frame=single,
	breaklines=true,
	breakatwhitespace=true,
	showstringspaces=false,
	tabsize=2,
	captionpos=b,
	literate=%
	{á}{{\'a}}1 {é}{{\'e}}1 {í}{{\'i}}1 {ó}{{\'o}}1 {ú}{{\'u}}1
	{Á}{{\'A}}1 {É}{{\'E}}1 {Í}{{\'I}}1 {Ó}{{\'O}}1 {Ú}{{\'U}}1
	{ñ}{{\~n}}1 {Ñ}{{\~N}}1 {¿}{{?`}}1 {¡}{{!`}}1
}

% ---------- Espacio reservado para imágenes faltantes ----------
\newcommand{\imgph}[2]{%
	\fbox{\begin{minipage}[c][#2][c]{0.92\columnwidth}
			\centering\textcolor{codegray}{\footnotesize [Espacio para imagen: #1]}
	\end{minipage}}%
}

\def\BibTeX{{\rm B\kern-.05em{\sc i\kern-.025em b}\kern-.08em
	T\kern-.1667em\lower.7ex\hbox{E}\kern-.125emX}}

\begin{document}

\title{Práctica 6: Módulos de Kernel y Dispositivos de Caracteres}

\author{
	\IEEEauthorblockN{Belmont Muñoz Samuel}
	\IEEEauthorblockA{\textit{Ingeniería en Computación} \\
		\textit{Facultad de Ingeniería}\\
		No. Cuenta: 317318289}
	\and
	\IEEEauthorblockN{Domínguez Palacios Jesús Alejandro}
	\IEEEauthorblockA{\textit{Ingeniería en Computación} \\
		\textit{Facultad de Ingeniería}\\
		No. Cuenta: 421066205}
	\and
	\IEEEauthorblockN{Dorantes Mejía Emmanuel}
	\IEEEauthorblockA{\textit{Ingeniería en Computación} \\
		\textit{Facultad de Ingeniería}\\
		No. Cuenta: 313004766}
	\and
	\IEEEauthorblockN{García Olvera Emiliano}
	\IEEEauthorblockA{\textit{Ingeniería en Computación} \\
		\textit{Facultad de Ingeniería}\\
		No. Cuenta: 314256009}
	\and
	\IEEEauthorblockN{Hernández Alonso Antonio Adonahí}
	\IEEEauthorblockA{\textit{Ingeniería en Computación} \\
		\textit{Facultad de Ingeniería}\\
		No. Cuenta: 320257780}
	\and
	\IEEEauthorblockN{Jiménez Reyes Alan Yael}
	\IEEEauthorblockA{\textit{Ingeniería en Computación} \\
		\textit{Facultad de Ingeniería}\\
		No. Cuenta: 320026726}
}

\maketitle

\begin{center}
	\vspace{-0.3cm}
	{\small
		\textbf{Práctica No.:} 6 \quad | \quad
		\textbf{Grupo:} 6 \quad | \quad
		\textbf{Fecha:} 12/10/2026
	}
	\vspace{0.2cm}
\end{center}

\begin{abstract}
	Este reporte documenta el desarrollo de módulos de kernel y drivers de dispositivos de caracteres en Linux embebido sobre Raspberry Pi 4. Se describe la compilación, carga y descarga de módulos mediante \texttt{insmod} y \texttt{rmmod}, el paso de parámetros a módulos y su inspección en \texttt{/sys/module}, y la implementación de un driver de caracteres virtual que expone las operaciones \texttt{open}, \texttt{read} y \texttt{write}. El documento aborda el uso de \texttt{copy\_from\_user} y \texttt{copy\_to\_user} para transferencia segura de datos entre espacio de kernel y espacio de usuario, la medición de costos de syscalls mediante \texttt{strace}, y la identificación de drivers utilizados por programas de prácticas anteriores. Se demuestra experimentalmente la frontera entre el espacio de usuario y el espacio de kernel del sistema operativo.
\end{abstract}

\begin{IEEEkeywords}
	Kernel, módulos, drivers, dispositivos de caracteres, insmod, rmmod, copy\_from\_user, strace, sysfs, Raspberry Pi
\end{IEEEkeywords}

\section{Objetivo}

\subsection{Objetivo General}
Compilar, cargar y descargar módulos de kernel en la Raspberry Pi 4, y construir un driver de dispositivos de caracteres virtual para comprender la frontera entre el espacio de usuario y el espacio de kernel.

\subsection{Objetivos Específicos}
\begin{itemize}
	\item Comprobar que los headers del kernel instalados corresponden al kernel en ejecución.
	\item Compilar un módulo mínimo, cargarlo con \texttt{insmod}, descargarlo con \texttt{rmmod}, y leer su salida con \texttt{dmesg}.
	\item Pasar parámetros a un módulo y consultarlos en \texttt{/sys/module}.
	\item Implementar las operaciones \texttt{open}, \texttt{read} y \texttt{write} de un dispositivo de caracteres utilizando \texttt{copy\_from\_user} y \texttt{copy\_to\_user}.
	\item Observar con \texttt{strace} las syscalls que llegan al driver y medir su costo frente a un acceso directo a memoria de usuario.
	\item Identificar, mediante \texttt{strace}, \texttt{lsmod} y \texttt{sysfs}, los archivos de dispositivo y los drivers utilizados por programas de prácticas anteriores (GPIO, I2C, SPI).
\end{itemize}

\section{Introducción}

El kernel de Linux es la capa de software que gestiona directamente el hardware, proporciona abstracciones para procesos de usuario, y garantiza el aislamiento y seguridad del sistema. Los módulos de kernel permiten extender la funcionalidad del kernel en tiempo de ejecución sin recompilar ni reiniciar el sistema completo.

\subsection{Espacio de Usuario vs. Espacio de Kernel}

Los sistemas operativos modernos dividen la memoria en dos regiones protegidas por hardware:

\begin{itemize}
	\item \textbf{Espacio de usuario (user space)}: Donde ejecutan las aplicaciones normales. Tiene acceso restringido a memoria, no puede ejecutar instrucciones privilegiadas, y no puede acceder directamente al hardware. Un error en espacio de usuario solo afecta al proceso que lo causó.

	\item \textbf{Espacio de kernel (kernel space)}: Donde ejecuta el kernel y sus módulos. Tiene acceso completo a toda la memoria, puede ejecutar cualquier instrucción del procesador, y controla directamente el hardware. Un error en espacio de kernel puede colapsar todo el sistema.
\end{itemize}

Las llamadas al sistema (syscalls) son el mecanismo por el cual el espacio de usuario solicita servicios del kernel, como abrir archivos, crear procesos, o acceder a dispositivos. La syscall causa una transición de modo usuario a modo kernel (privileged mode), ejecuta la operación solicitada, y retorna al modo usuario con el resultado.

\subsection{Módulos de Kernel}

Un módulo de kernel es código que puede cargarse y descargarse dinámicamente del kernel en ejecución. Los módulos típicamente implementan drivers de dispositivos, sistemas de archivos, o protocolos de red. Las ventajas de los módulos incluyen:

\begin{itemize}
	\item No requieren recompilar el kernel completo.
	\item Pueden cargarse solo cuando se necesitan, reduciendo el uso de memoria.
	\item Permiten desarrollo y depuración sin reiniciar el sistema.
	\item Facilitan la distribución de drivers propietarios sin modificar el código del kernel.
\end{itemize}

Los módulos se compilan contra los headers del kernel instalado, que deben coincidir exactamente con la versión del kernel en ejecución. El comando \texttt{uname -r} muestra la versión actual, y los headers se instalan mediante \texttt{apt install raspberrypi-kernel-headers}.

\subsection{Dispositivos de Caracteres}

Linux clasifica los dispositivos en tres categorías principales: dispositivos de bloque (discos), dispositivos de caracteres (terminales, puertos serie, sensores), y dispositivos de red. Los dispositivos de caracteres son aquellos que se acceden como flujos de bytes sin estructura de bloques definida.

Un driver de caracteres implementa un conjunto de operaciones definidas en la estructura \texttt{file\_operations}:

\begin{lstlisting}[language=C]
struct file_operations {
    int (*open)(struct inode *, struct file *);
    int (*release)(struct inode *, struct file *);
    ssize_t (*read)(struct file *, char __user *, size_t, loff_t *);
    ssize_t (*write)(struct file *, const char __user *, size_t, loff_t *);
    // ... otras operaciones
};
\end{lstlisting}

El kernel invoca estas funciones cuando el espacio de usuario ejecuta syscalls como \texttt{open()}, \texttt{read()}, o \texttt{write()} sobre el archivo de dispositivo correspondiente.

\subsection{Transferencia Segura de Datos}

El kernel no puede acceder directamente a punteros del espacio de usuario debido a razones de seguridad y protección de memoria. Las funciones \texttt{copy\_from\_user()} y \texttt{copy\_to\_user()} realizan transferencias verificadas:

\begin{lstlisting}[language=C]
// Copia desde usuario hacia kernel
unsigned long copy_from_user(void *to, const void __user *from, unsigned long n);

// Copia desde kernel hacia usuario
unsigned long copy_to_user(void __user *to, const void *from, unsigned long n);
\end{lstlisting}

Estas funciones verifican que las direcciones de memoria sean válidas, que el proceso tenga permisos de acceso, y manejan fallos de página si la memoria está paginada a disco. Retornan el número de bytes que no pudieron copiarse (0 indica éxito completo).

\subsection{Herramientas de Inspección}

\begin{itemize}
	\item \texttt{lsmod}: Lista módulos cargados actualmente en el kernel.
	\item \texttt{modinfo}: Muestra información sobre un módulo (parámetros, autor, licencia).
	\item \texttt{dmesg}: Lee el buffer circular del kernel donde se registran mensajes mediante \texttt{printk()}.
	\item \texttt{strace}: Intercepta y registra todas las syscalls ejecutadas por un proceso, mostrando argumentos y resultados.
	\item \texttt{/sys/module/}: Sistema de archivos virtual que expone información sobre módulos cargados, incluyendo parámetros.
\end{itemize}

\subsection{Conceptos Aplicados en Esta Práctica}

La Tabla~\ref{tab:conceptos} resume los conceptos fundamentales de kernel y su aplicación concreta en las distintas partes de esta práctica. Cada concepto que antes era invisible se hará explícito mediante herramientas de inspección y módulos creados manualmente.

\begin{table}[H]
	\centering
	\caption{Conceptos de kernel aplicados en esta práctica}
	\label{tab:conceptos}
	\small
	\begin{tabular}{|p{2.2cm}|p{3.2cm}|p{2.2cm}|}
		\hline
		\textbf{Concepto} & \textbf{Definición} & \textbf{En esta práctica} \\
		\hline
		Espacio de usuario & Región donde corren los programas comunes. Tiene acceso restringido: no puede tocar el hardware ni la memoria de otros procesos. & \texttt{cliente.c} y \texttt{tiempo.c} (Partes E y F), igual que todo el código de las prácticas 1 a 5. \\
		\hline
		Espacio de kernel & Región donde corre el kernel de Linux. Tiene acceso total al hardware y a la memoria de todos los procesos. & \texttt{hola.ko}, \texttt{parametros.ko} y \texttt{p6buf.ko}. \\
		\hline
		Llamada al sistema (syscall) & Mecanismo con el que un programa de usuario pide un servicio al kernel: \texttt{open}, \texttt{read}, \texttt{write}, \texttt{ioctl}, \texttt{pipe}. & Las que \texttt{strace} muestra al ejecutar cliente (Parte E). \\
		\hline
		Cambio de modo & Al hacer una syscall el procesador guarda el contexto del programa, pasa a modo kernel, ejecuta el servicio y regresa. Cada cruce consume tiempo. & Se mide en la Parte F. \\
		\hline
		Módulo del kernel & Archivo \texttt{.ko} con código de kernel que se carga y se descarga sin reiniciar el sistema. & \texttt{insmod}, \texttt{lsmod} y \texttt{rmmod} (Partes B a D). \\
		\hline
		Driver & Código de kernel específico de un dispositivo. Traduce las syscalls genéricas en las operaciones que ese dispositivo entiende. Puede ser un módulo o venir integrado en la imagen del kernel. & \texttt{p6buf} (Parte D) y los drivers reales de la Parte G. \\
		\hline
		Archivo de dispositivo & Entrada de \texttt{/dev} por la que un programa de usuario llega a un driver. & \texttt{/dev/p6buf} y \texttt{/dev/gpiochip0}. \\
		\hline
	\end{tabular}
\end{table}

\subsection{Relación con Prácticas Anteriores}

En prácticas anteriores se utilizaron drivers del kernel de forma transparente: \texttt{gpiod} para control de GPIO (Práctica 2), \texttt{i2c-dev} para comunicación I2C, y \texttt{spidev} para comunicación SPI (Práctica 3). Esta práctica desmitifica esos componentes, mostrando cómo están implementados y cómo interactúan con el espacio de usuario.

\section{Desarrollo}

\subsection{Material y Equipo}

Para la realización de esta práctica se utilizaron los siguientes elementos:

\begin{itemize}
	\item Raspberry Pi 4 Model B con Raspberry Pi OS instalado
	\item Headers del kernel (\texttt{raspberrypi-kernel-headers}) correspondientes a la versión en ejecución
	\item Toolchain de desarrollo en C (GCC) con soporte para compilación de módulos
	\item Repositorio Git del curso con programas de prácticas anteriores
	\item Laptop con acceso SSH al sistema (o teclado y monitor conectados directamente)
	\item Conexión a Internet para instalación de headers del kernel
	\item Circuito de LEDs/sensores de prácticas 2 y 3 (opcional, no requerido para funcionamiento básico)
\end{itemize}

\textbf{Nota}: No se requiere hardware adicional. El driver desarrollado es virtual y no controla dispositivos físicos reales, aunque los programas de prácticas anteriores que utilizan GPIO, I2C y SPI se analizarán para identificar sus interacciones con drivers del kernel.

\section{Resultados y Análisis}

Contenido pendiente: Presentar los resultados obtenidos de cada parte del desarrollo.

\section{Conclusiones}

\textbf{Belmont Muñoz Samuel}. \\
Conclusión pendiente.

\textbf{Domínguez Palacios Jesús Alejandro}. \\
Conclusión pendiente.

\textbf{Dorantes Mejía Emmanuel}. \\
Conclusión pendiente.

\textbf{García Olvera Emiliano}. \\
Conclusión pendiente.

\textbf{Hernández Alonso Antonio Adonahí}. \\
Conclusión pendiente.

\textbf{Jiménez Reyes Alan Yael}. \\
Conclusión pendiente.

\begin{thebibliography}{00}
	\bibitem{lkd} Robert Love. (2010). \textit{Linux Kernel Development}, 3rd ed. Addison-Wesley.
	\bibitem{ldd3} Jonathan Corbet, Alessandro Rubini, Greg Kroah-Hartman. (2005). \textit{Linux Device Drivers}, 3rd ed. O'Reilly Media. Disponible en: https://lwn.net/Kernel/LDD3/
	\bibitem{kernel_docs} The kernel development community. (2024). ``The Linux Kernel Documentation''. Disponible en: https://www.kernel.org/doc/html/latest/
	\bibitem{module_programming} Peter Jay Salzman et al. (2024). ``The Linux Kernel Module Programming Guide''. Disponible en: https://sysprog21.github.io/lkmpg/
	\bibitem{claude_code} Anthropic. (2025). ``Claude Code: AI-Powered Development Assistant''. Disponible en: https://claude.com/claude-code
	\bibitem{strace_man} strace developers. (2024). ``strace(1) - Linux man page''. Disponible en: https://man7.org/linux/man-pages/man1/strace.1.html
\end{thebibliography}

\end{document}
